\begin{frame}[label=q5a1]{5a.1 · DataFrame: agrupar y filtrar dos columnas}
  \small
  \pregunta{¿Qué combinaciones de departamento y rango de edad reúnen al menos 1\,000 personas reportadas?}
  \entrada{93\,831 filas. \texttt{departamento}, \texttt{rango\_edad}, \texttt{cantidad}.}
  \flujo{93\,831 registros}{\texttt{groupBy}\\Departamento y edad}{\texttt{SUM(cantidad)}\\\texttt{filter >= 1000}}{30 combinaciones\\Orden descendente}
  \begin{columns}[T,onlytextwidth]
    \column{0.55\textwidth}
    \textbf{Código Scala original (fragmento)}
    \codigo{sql/05a_1_departamento_rangoedad}{10}{16}
    \column{0.42\textwidth}
    \textbf{Salida}\\[5pt]
    {\footnotesize
\begin{tabular}{@{}lll@{}}
\toprule
Departamento & Edad & Personas\\
\midrule
Lima Met. & 12--17 & 17\,901\\
Lima Met. & 18--29 & 7\,875\\
Junín & 12--17 & 4\,169\\
\bottomrule
\end{tabular}
    }
  \end{columns}
  \respuesta{Lima Metropolitana y el rango 12--17 encabezan el resultado. El filtro opera sobre el total agregado, no sobre la cantidad de cada fila.}
  \origen{sql}{05a_1_departamento_rangoedad}
  \note{25 s. Señalar las dos claves del groupBy, luego el umbral. Los otros tres ejemplos 5a están enlazados en el mapa.}
\end{frame}
